\section{Experimental Design}

\subsection{Research Questions}

We organize the evaluation around five questions:

\begin{description}
\item[RQ1: Fixed-budget effectiveness.] Does \system{} solve more tasks than
    complete-system baselines under the same component-wise hard resource
    vector?
    \item[RQ2: Efficiency.] At noninferior accepted rate, does it reduce
    unconditional token, checker, latency, or monetary expenditure per task?
    \item[RQ3: Mechanism attribution.] Holding the structured harness fixed,
    what is the incremental effect of the expanded action mask, value-per-cost
    selection, and remaining-budget admission?
    \item[RQ4: Model routing.] Holding the action space and selector fixed,
    does adaptive cheap/strong routing outperform a preregistered fixed
    escalation schedule with the same model portfolio and caps?
    \item[RQ5: Task regime and overhead.] When does the end-to-end benefit of
    structured control outweigh its projection, prompting, and controller
    overhead relative to minimal refinement?
\end{description}

\subsection{Tasks and Data Separation}

The confirmatory evaluation starts from fixed Lean~4 statements, imports, and
proof holes. This isolates proof-search allocation from autoformalization. We
use two public benchmark test splits frozen before any test-arm execution:

\begin{enumerate}
    \item miniF2F Olympiad-style statements
    \cite{zheng2021minif2f}; the current adapter freezes 244 validation and
    244 test tasks and verifies all 488 masked scaffolds with the benchmark's
    real Lean toolchain;
    \item \todo{freeze the second public Lean benchmark, immutable revision,
    test split, toolchain, and eligible task count}.
\end{enumerate}

For every suite we publish the source revision, license, Lean and Mathlib
revisions, split, eligible task count, task identifiers, and statement hashes.
Validation tasks are used only to choose budget vectors, fixed-escalation
thresholds, score ranks or weights, and frozen cost-history snapshots. Test
tasks are not used to remove tasks, tune the controller, or update cost and
success estimates. Ground-truth proofs are unavailable to generation,
retrieval, prompts, and failure context.

A local multi-obligation suite is a secondary mechanistic stress test for
dependency unlocks, helper reuse, and final assembly. It is reported separately
and is not pooled into the confirmatory public-benchmark result. Natural-language
formalization is outside the primary causal comparison because it introduces a
distinct failure process and resource-allocation problem.

\subsection{End-to-End Baselines}

Complete-system comparisons include:

\begin{enumerate}
    \item \textbf{Independent sampling}: independent whole-proof candidates
    with no feedback reuse;
    \item \textbf{Direct strong}: the strong-tier boundary policy under the
    same resource caps;
    \item \textbf{Minimal refinement}: propose--check--remember--refine with
    compact memory and the same retrieval corpus, checker, safety review, and
    model access policy;
    \item \textbf{Fixed escalation}: minimal refinement with a
    validation-frozen cheap-to-strong schedule; and
    \item \system: the full structured action, selector, admission, and
    routing policy.
\end{enumerate}

Where code and task compatibility permit, we additionally compare with the
continue-versus-redecompose cost-aware router
\cite{rognvaldsson2026costquality}. Otherwise it remains a conceptual
comparison and is not assigned a non-equivalent numerical result.

These baselines answer whether the complete controller pays for its overhead.
They are \emph{not} used to attribute a difference to state representation,
action-space expansion, or selection in isolation, because the systems have
different internal contexts and control loops.

\subsection{Controlled Mechanism Comparisons}

Mechanism attribution uses one shared structured harness. Every controlled arm
has the same workspace, context projection, proposer template, retrieval
corpus, checker, safety review, decoding setup, paired seed, cost-history
snapshot, final-check reserve, and component-wise hard budget. Only the named
factor changes.

\begin{table}[ht]
\centering
\small
\begin{tabular}{llll}
\toprule
Arm & Action mask & Selector & Model policy \\
\midrule
\(A_{00}\) & implement/repair & value only & one common model \\
\(A_{10}\) & expanded & value only & one common model \\
\(A_{01}\) & implement/repair & full cost-aware & one common model \\
\(A_{11}\) & expanded & full cost-aware & one common model \\
\(S_{\mathrm{ratio}}\) & expanded & value/cost, no pre-admission & one common model \\
\(R_{\mathrm{fixed}}\) & expanded & full cost-aware & fixed tier schedule \\
\system & expanded & full cost-aware & adaptive tier routing \\
\bottomrule
\end{tabular}
\caption{Controlled arms. The expanded mask adds capability tests,
decomposition, argument refinement, and representation change to implementation
and repair. Retrieval access inside proof calls is held fixed. The full
cost-aware selector adds remaining-vector admission to value-per-cost ranking.}
\label{tab:controlled-arms}
\end{table}

The \(2\times2\) cells \(A_{00},A_{10},A_{01},A_{11}\) estimate action-space
and selector effects under a single common model, including their interaction.
Within the expanded mask, \(A_{10}\rightarrow S_{\mathrm{ratio}}\) isolates
cost normalization and \(S_{\mathrm{ratio}}\rightarrow A_{11}\) isolates
estimated remaining-budget admission. The contrast
\(R_{\mathrm{fixed}}\rightarrow\system\) isolates adaptive routing because both
arms expose the identical cheap/strong endpoints, total request limit,
strong-tier cap, token and monetary caps, selector, and action mask. Cheap-only
and strong-only runs are reported as boundary policies, not routing ablations.

We also include an attempt/terminate mask as a weak action-space reference when
the common proposer protocol supports it. Component removals for unlock value,
hypothesis discrimination, diagnostic repair, capability audits, and
final-check reserves are exploratory: they can alter both state visitation and
action availability. Controlled marginal effects are not assumed to add up to
the full-system difference; interaction terms are reported directly.

\subsection{Matched Resource Budgets}

We preregister a nested sequence of component-wise budget profiles,
\[
  \mathbf B^{(1)}\preceq\mathbf B^{(2)}\preceq\cdots
  \preceq\mathbf B^{(K)},
\]
where a profile limits model requests, logical and provider-reported tokens,
checker calls, wall time, standardized monetary cost, and tier-specific calls.
\todo{Freeze the numeric low, medium, and high profiles from validation-only
pilots.} All arms in one benchmark--model--budget cell receive the same vector,
timeout, final-assembly accounting, and frozen price table. An arm may stop
early; it is never forced to consume its full allowance.

Each budget profile is run independently. A low-budget result is not obtained
by truncating a high-budget trace because the controller observes remaining
resources and may change its decisions. If wall time is used as a confirmatory
constraint, hardware, concurrency, server warm-up, and task ordering are fixed;
otherwise latency is descriptive. Independent sampling is run sequentially
unless concurrency itself is explicitly matched and reported.

Projection and prompt tokens, proposal generation, retrieval, every checker
call, routing, and final assembly are charged. We report two monetary views:
standardized cost computed from the frozen price table, and actual
provider-billed cost. Cached input, hidden reasoning-token availability, and
billed-cost measurement coverage are shown separately. Missing measurements
are never imputed as zero.

\subsection{Outcomes, Metrics, and Estimands}

The primary endpoint is \emph{accepted under \(\mathbf B\)}: the run must
produce a fully integrated proof accepted by both the Lean checker and safety
review; multi-obligation tasks additionally require final assembly to pass.
The accepted-rate denominator is accepted plus proof failure. Provider,
transport, checker-server, and external timeout failures are classified as
infrastructure, retried under the same frozen configuration, excluded from the
primary denominator if unresolved, and reported with selected, completed,
retried, missing, and per-arm counts.

For each resource \(j\), we distinguish:

\begin{itemize}
    \item unconditional mean \(\mathbb E[C_j]\) per attempted task, including
    proof failures---the primary efficiency quantity;
    \item successful-run cost \(\mathbb E[C_j\mid Y=1]\), a descriptive
    statistic that may change with the set of tasks solved; and
    \item amortized cost per accepted proof
    \(\mathbb E[C_j]/P(Y=1)\).
\end{itemize}

We report accepted rate and actual resource use at every preregistered vector.
Success--cost curves sweep one declared cap or a nested budget sequence while
holding the other definitions fixed; any area-under-curve statistic uses a
preregistered normalized interval. Additional metrics include restricted mean
cost to solution with failures treated as censored at the cap, strong-tier
calls and conditional outcomes, cap-hit frequencies, projection/control
overhead, and checks outside the ancestry of the final accepted artifact.

Mechanism evidence includes, for every richer-only action kind, the number of
eligible states, proposals, executions, positive verifier-grounded transitions,
accepted artifacts, and exclusive paired successes. Immediate success
probabilities are assessed with reliability plots and Brier score when a
probabilistic model exists. A non-probabilistic priority \(Q\) is evaluated by
ranking lift and outcome by score quantile, not called calibrated.

\subsection{Statistical Analysis and Claim Criteria}

Task and repetition identifiers are paired across arms, and arm order is
rotated or randomized. Confidence intervals use a paired hierarchical
bootstrap with task as the resampling cluster; repeated generations for the
same task are not treated as independent tasks. Accepted-rate differences are
reported in percentage points with paired intervals for each
benchmark--model--budget cell. Overall results are secondary and state whether
they are macro-averaged across cells or weighted by task count.

Before test execution we freeze the number of repetitions and an accepted-rate
non-inferiority margin of \todo{1 percentage point, or another justified value
chosen before test runs}. The allocation claim is supported if \system{} either
(i) improves accepted rate at the same resource vector, or (ii) is noninferior
in accepted rate while reducing unconditional resource use, with the
corresponding paired interval. A gain of the full system over minimal refinement
without a gain over the matched structured selector control is reported as an
action-space or representation result, not as evidence for cost-aware
allocation.

\subsection{Reproducibility and Leakage Controls}

Each run freezes benchmark provenance, ordered task identifiers, model
versions, temperature and seed, prompt/tool protocol, action mask, selector,
routing thresholds, price table, history snapshot, checker configuration, and
all budget coordinates behind one configuration hash. Resume is rejected when
the hash changes. Traces record raw candidates, checker categories, goal
fingerprints, branch and obligation versions, complete selector choice sets,
model and token use, elapsed time, stop reason, estimate errors, measurement
coverage, and the final workspace. Summary tables are regenerated from
per-task results and append-only traces rather than edited by hand.
